\section{From size to independence}\label{sec:size}

We next show that a vertex set larger than both $k$ and twice the clique
number contains three independent vertices.  The argument uses only the
forbidden cycle and the clique bound.  We first prove two elementary facts
about graphs with independence number at most two.
The first is a special case of the Chv\'atal--Erd\H{o}s Hamiltonicity
theorem~\cite[Theorem~1]{ChvatalErdos1972}. The shortening argument in the
second appears in Radziszowski and Jin's proof of their pancyclicity
theorem~\cite[Theorem~2]{RadziszowskiJin1994}. We include both proofs.

\begin{lemma}\label{size:hamiltonian}
Every finite simple 2-connected graph $H$ with $\alpha(H)\le2$ contains a
cycle through all its vertices.
\end{lemma}

\begin{proof}
A 2-connected graph has minimum degree at least two and therefore contains
a cycle: all neighbors of an endpoint of a longest path lie on that path,
and two such neighbors give a cycle.  Choose a cycle $C$ of maximum order,
orient it, and write $a^+$ for the successor of $a\in V(C)$.

Suppose that a component $D$ of $H-V(C)$ exists.  Its set of neighbors on
$C$ contains at least two vertices.  Indeed, it is nonempty by
connectedness, and a unique such vertex would be a cutvertex of $H$.
If $a\in V(C)$ has a neighbor in $D$, then $a^+$ has none.  Otherwise a
path from $a$ to $a^+$ with nonempty interior in the connected graph $D$
could replace the edge $aa^+$, giving a longer cycle.  This path exists
even if $a$ and $a^+$ have the same neighbor in $D$.

Choose distinct vertices $a,b\in V(C)$ with neighbors in $D$, and a path
$P$ from $a$ to $b$ with nonempty interior in $D$.  The vertices $a^+,b^+$
are distinct, neither equals $a$ or $b$, and both have no neighbor in $D$.
If $a^+b^+$ were an edge, we could traverse $P$ from $a$ to $b$, follow
$C$ backwards from $b$ to $a^+$, use $a^+b^+$, and follow $C$ forwards
from $b^+$ to $a$.  Deleting $aa^+$ and $bb^+$ from $C$ leaves exactly
the two cycle segments used here.  Hence this would be a cycle containing
every vertex of $C$ and at least one additional vertex, a contradiction.
Thus $a^+$ and $b^+$ are nonadjacent.  Together with any vertex of $D$,
they form an independent triple, also a contradiction.  Therefore $C$
is spanning.
\end{proof}

\begin{lemma}\label{size:shortening}
Let $H$ be a finite simple graph on $r\ge7$ vertices with
$\alpha(H)\le2$.  If $H$ has a Hamiltonian cycle, then it has a cycle
on exactly $r-1$ vertices.
\end{lemma}

\begin{proof}
Label the vertices around a Hamiltonian cycle by $0,1,\ldots,r-1$, with
indices interpreted modulo $r$.  If some edge joins $i$ to $i+2$, it
bypasses the vertex $i+1$ and gives the required cycle.  Suppose there
is no such edge.  In each triple $i,i+2,i+4$, the pairs $i(i+2)$ and
$(i+2)(i+4)$ are nonedges.  Since the triple is not independent,
$i(i+4)$ is an edge.  All cyclic step-four edges are therefore present.

Consider the cyclic sequence
\begin{equation}\label{size:shortening-cycle}
 1,2,\ldots,r-6,\quad r-2,r-1,\quad r-5,r-4,r-3,\quad1.
\end{equation}
Before closing, its three blocks partition the vertices $1,\ldots,r-1$.
Consecutive vertices within a block are adjacent on the original cycle.
The three edges joining the blocks and closing the sequence are
\[
 (r-6)(r-2),\qquad (r-1)(r-5),\qquad (r-3)1;
\]
their cyclic differences are $4,-4,4$, respectively.  Thus
\eqref{size:shortening-cycle} is a cycle of order $r-1$.
When $r=7$, the first block is the singleton $1$, and the sequence is
explicitly $1,5,6,2,3,4,1$, so the same construction applies.
\end{proof}

\begin{lemma}[Size test]\label{size:test}
Let $k\ge5$ and $t\ge1$ be integers.  Let $G$ be a finite simple graph
with no cycle on exactly $k+1$ vertices and with clique number at most
$t$.  For every $Y\subseteq V(G)$,
\[
 |Y|>\max\{k,2t\}\quad\Longrightarrow\quad\alpha(G[Y])\ge3.
\]
\end{lemma}

\begin{proof}
Suppose that $H=G[Y]$ has $\alpha(H)\le2$ and order
$n>\max\{k,2t\}$.  If $H$ is disconnected, it has exactly two
components, each a clique.  Indeed, three components give an independent
triple, as does a nonadjacent pair in one component together with a vertex
of another.  The two cliques have at most $t$ vertices each, contrary to
$n>2t$.

If $H$ has a cutvertex $v$, the same argument shows that $H-v$ has
exactly two components $A,B$, both cliques of order at most $t$.
The inequalities
\[
 |A|+|B|+1=n>2t,\qquad |A|,|B|\le t
\]
force $|A|=|B|=t$.  The vertex $v$ must be adjacent to all of one
component: a nonneighbor in each would form an independent triple with
$v$.  That component together with $v$ is a clique of order $t+1$,
again a contradiction.

Consequently $H$ is 2-connected and has a Hamiltonian cycle by
Lemma~\ref{size:hamiltonian}.  Its order is $n\ge k+1$.  If a current
cycle has order $r>k+1$, the graph induced by its vertices is Hamiltonian,
has independence number at most two, and satisfies $r\ge k+2\ge7$.
Lemma~\ref{size:shortening} therefore produces a cycle of order $r-1$.
Repeated shortening reaches exactly $k+1$, contradicting the hypothesis
on $G$.
\end{proof}
